\documentclass[11pt,a4paper]{article}

\usepackage[dutch]{babel}
\usepackage[table]{xcolor}
\usepackage{longtable}
\usepackage{booktabs}
\usepackage{array}
\usepackage{pdflscape}
% \usepackage{geometry}

\renewcommand{\arraystretch}{1.3}

% Kolomtypes (m{..} = verticaal midden in de cel)
\newcolumntype{C}[1]{>{\centering\arraybackslash}m{#1}}
\newcolumntype{L}[1]{>{\raggedright\arraybackslash}m{#1}}

\newcommand{\XX}{\bfseries\large X}

\usepackage[a4paper,
    left=1.0cm,
    right=1.0cm,
    top=0cm,
    bottom=0cm]{geometry}

\begin{document}
\pagestyle{empty}

\begin{landscape}
\small

\begin{longtable}{
L{2.cm}    % Beroepsvaardigheid
L{5cm}      % Doel
L{5cm}||     % Handelingsindicatoren
C{0.7cm}|     % PVA
C{0.9cm}|   % Eindverslag
C{0.8cm}|     % Poster
C{1.5cm}||   % Poster-presentatie
C{1.2cm}|   % Stagebezoek week 5
C{1.2cm}|   % SLBgesprek + zelfreflectie
C{1.7cm}|   % Praktijk
C{1.7cm}||      % Logboek
C{1.3cm}
}

\textbf{Beroeps-vaardigheid} &
\textbf{Doel voor dit jaar} &
\textbf{Handelingsindicatoren} &
\textbf{\scriptsize
 PVA} &
\textbf{\scriptsize
Eind\-verslag} &
\textbf{\scriptsize
Poster} &
\textbf{\scriptsize
Poster\-presentatie} &
\textbf{\scriptsize
Feedback week 5} &
\textbf{\scriptsize
Feedback week 15} &
\textbf{\scriptsize
Labjournaal \& logboek} &
\textbf{\scriptsize
Criterium Gericht Interview} &
\textbf{\scriptsize
SLB-gesprek}  \\
\midrule
\endfirsthead

\toprule
\textbf{Beroeps-vaardigheid} &
\textbf{Doel voor dit jaar} &
\textbf{Handelingsindicatoren} &
\textbf{\scriptsize
 PVA} &
\textbf{\scriptsize
Eind\-verslag} &
\textbf{\scriptsize
Poster} &
\textbf{\scriptsize
Poster\-presentatie} &
\textbf{\scriptsize
Feedback week 5} &
\textbf{\scriptsize
Feedback week 15} &
\textbf{\scriptsize
Labjournaal / logboek} &
\textbf{\scriptsize
Criterium Gericht Interview} &
\textbf{\scriptsize
SLB-gesprek} \\
\midrule
\endhead

%====================================================
\rowcolor{gray!20}
\multicolumn{3}{l}{\textbf{ONDERZOEKSVAARDIGHEDEN}} &
\multicolumn{4}{>{\cellcolor{green!80!black}}l}{} &
\multicolumn{4}{>{\cellcolor{orange!90}}l}{} &
\multicolumn{1}{>{\cellcolor{blue!60}}l}{} \\
\midrule

Analytische vaardigheden &
Je kan vragen en deelvragen van een onderzoek formuleren en relevante bronnen gebruiken. &
In samenspraak een probleem analyseren en vertalen naar doelstellingen van de onderzoeksopdracht. &
\XX & & & \XX & \XX & & & & \\

\midrule
Modelleer vaardigheden &
Je kan een geschikt model selecteren, opstellen en aanpassen aan de situatie. &
Uit aangereikte modellen de juiste selecteren, aannames formuleren en het model valideren. &
\XX & \XX & &\XX  & \XX & \XX & \XX& \XX &\\

\midrule
Experimentele vaardigheden &
Je kan een experimentele aanpak opstellen en resultaten zorgvuldig registreren. &
Een geschikte methode kiezen, experiment uitvoeren en reproduceerbare resultaten documenteren. &
\XX & \XX & &\XX  & \XX & \XX & \XX & \XX &\\

%====================================================
\rowcolor{gray!20}
\multicolumn{12}{l}{\textbf{PROJECTVAARDIGHEDEN}} \\
\midrule

Project opzetten &
Projectplan maken met planning en taakverdeling. &
Werkplan opstellen rekening houdend met randvoorwaarden en PvE. &
\XX & & & & \XX& \XX & \XX& \XX &\XX\\

\midrule
Project uitvoeren &
Project volgens plan uitvoeren en waar nodig bijstellen. &
Voortgang monitoren, knelpunten signaleren en communiceren. &
 & & & & \XX& \XX & \XX & \XX & \XX\\

\midrule
Project opleveren &
Resultaten samenvatten en vervolgstappen formuleren. &
Conclusies trekken en documentatie opleveren volgens richtlijnen. &
 & \XX & \XX & \XX & & & \XX& &\XX\\

%====================================================
\rowcolor{gray!20}
\multicolumn{12}{l}{\textbf{DATAVAARDIGHEDEN}} \\
\midrule

Data verzamelen en bewerken &
Data genereren en bewerken. &
Meetresultaten verwerken en beoordelen op realisme. &
 & \XX & & & & \XX& \XX& \XX &\\

\midrule
Data analyse en statistiek &
Statistische analyse uitvoeren. &
Betrouwbaarheid beoordelen en validatie uitvoeren. &
 & \XX & \XX & \XX & & & &  &\\

\midrule
Data visualiseren en interpreteren &
Data visualiseren en interpreteren. &
Resultaten grafisch weergeven en fysisch interpreteren. &
 & \XX & \XX & \XX & & & \XX& &\\

\newpage
%====================================================
\rowcolor{gray!20}
\multicolumn{3}{l}{\textbf{PERSOONLIJKE VAARDIGHEDEN}} &
\multicolumn{4}{>{\cellcolor{green!80!black}}l}{} &
\multicolumn{4}{>{\cellcolor{orange!90}}l}{} &
\multicolumn{1}{>{\cellcolor{blue!60}}l}{} \\
\midrule

Reflectie en feedback &
Leerdoelen formuleren en reflecteren. &
Feedback vragen en verwerken. &
 & & & & \XX & \XX & &  &\XX\\

\midrule
Interpersoonlijke vaardigheden &
Actief functioneren in een team. &
Communiceren met opdrachtgever en teamleden. &
 & & & & \XX & \XX & & & \XX\\

\midrule
Loopbaan-ontwikkeling &
Oriëntatie op professionele ontwikkeling. &
Bewust zijn van professionele en ethische aspecten. &
 & & & & \XX & \XX & & & \XX\\

%====================================================
\rowcolor{gray!20}
\multicolumn{12}{l}{\textbf{IMPACT}} \\
\midrule

Boodschap structureren &
Onderzoeksresultaten logisch presenteren. &
Resultaten structureren en koppelen aan onderzoeksvraag. &
 & \XX & \XX & & & & &  &\\

\midrule
Mondeling presenteren &
Gestructureerde presentatie geven. &
Duidelijk en passend bij publiek presenteren. &
 & & & \XX & & & & & \\

\midrule
Schriftelijk rapporteren &
Rapport volgens professionele standaard opleveren. &
Resultaten integreren in één consistent document. &
 \XX & \XX & & & & & &  &\\

\bottomrule
\end{longtable}

\end{landscape}
\end{document}
